\documentclass[letterpaper]{article} 
\usepackage[utf8]{inputenc}
\linespread{0.85}
\usepackage[T1]{fontenc}
\usepackage{amsmath}
\usepackage{amsfonts}
\usepackage{amssymb}
\usepackage{array}
\usepackage{fontspec}
\usepackage{booktabs}
\usepackage{hyperref}
\usepackage[version=4]{mhchem}
\usepackage{stmaryrd}
\usepackage[dvipsnames]{xcolor}
\colorlet{LightRubineRed}{RubineRed!70}
\colorlet{Mycolor1}{green!10!orange}
\definecolor{Mycolor2}{HTML}{00F9DE}
\usepackage{graphicx}
\usepackage{amsmath}
\usepackage{graphicx}
\usepackage{capt-of}
\usepackage{lipsum}
\usepackage{fancyvrb}
\usepackage{tabularx}
\usepackage{listings}
\usepackage[export]{adjustbox}
\graphicspath{ {./images/} }
\usepackage[utf8]{inputenc}
\usepackage[english]{babel}
\usepackage{float}
\usepackage{lipsum}
\usepackage{graphicx}
\usepackage{float}
\usepackage[margin=0.7in]{geometry}
\usepackage{amsmath}
\usepackage{graphicx}
\usepackage{capt-of}
\usepackage{tcolorbox}
\usepackage{lipsum}
\usepackage{graphicx}
\usepackage{float}
\usepackage{listings}
\definecolor{deepred}{RGB}{139,0,0}
\usepackage{hyperref} 
\usepackage{xcolor} % For custom colors
\lstset{
	language=Python,                % Choose the language (e.g., Python, C, R)
	basicstyle=\ttfamily\small, % Font size and type
	keywordstyle=\color{blue},  % Keywords color
	commentstyle=\color{gray},  % Comments color
	stringstyle=\color{red},    % String color
	numbers=left,               % Line numbers
	numberstyle=\tiny\color{gray}, % Line number style
	stepnumber=1,               % Numbering step
	breaklines=true,            % Auto line break
	backgroundcolor=\color{black!5}, % Light gray background
	frame=single,               % Frame around the code
}
\usepackage{float}
\usepackage[]{amsthm} %lets us use \begin{proof}
	\usepackage[]{amssymb} %gives us the character \varnothing
	
	
	\title{Homework 3, STAT 5291}
	\author{Zongyi Liu}
	\date{Sun, Mar 22, 2026}
	\begin{document}
		\maketitle
		
		\section{Question 1}
		{\color{deepred}\fontspec{Georgia}  Two-Way Factorial Application}
		
		This problem is an adaptation from the textbook \textit{Statistical Methods for Psychology}, by David C. Howell. Consider a study which includes the factors \textit{Age}, as well as \textit{Recall Condition}. The study included 50 participants in the 18-to-30-year age range, as well as 50 participants in the 55-to-65-year age range. The response variable is the number items the respondent recalled, for example, how many images the respondent recalled. The data in Table 1 have been created to have the same means and standard deviations as those reported in the original experiment. Assume that the data themselves are approximately normally distributed with acceptably equal variances. Also note that the data can be retrieved from Canvas (\texttt{HW4Problem1.csv}).
		
		\begin{table}[h]
			\centering
			\begin{tabular}{|l|ccccc|}
				\hline
				& Counting & Rhyming & Adjective & Imagery & Intention \\
				\hline
				Old & 9  & 7  & 11 & 12 & 10 \\
				Old & 8  & 9  & 13 & 11 & 19 \\
				Old & 6  & 6  & 8  & 16 & 14 \\
				Old & 8  & 6  & 6  & 11 & 5  \\
				Old & 10 & 6  & 14 & 9  & 10 \\
				Old & 4  & 11 & 11 & 23 & 11 \\
				Old & 6  & 6  & 13 & 12 & 14 \\
				Old & 5  & 3  & 13 & 10 & 15 \\
				Old & 7  & 8  & 10 & 19 & 11 \\
				Old & 7  & 7  & 11 & 11 & 11 \\
				\hline
				Young & 8  & 10 & 14 & 20 & 21 \\
				Young & 6  & 7  & 11 & 16 & 19 \\
				Young & 4  & 8  & 18 & 16 & 17 \\
				Young & 6  & 10 & 14 & 15 & 15 \\
				Young & 7  & 4  & 13 & 18 & 22 \\
				Young & 6  & 7  & 22 & 16 & 16 \\
				Young & 5  & 10 & 17 & 20 & 22 \\
				Young & 7  & 6  & 16 & 22 & 22 \\
				Young & 9  & 7  & 12 & 14 & 18 \\
				Young & 7  & 7  & 11 & 19 & 21 \\
				\hline
			\end{tabular}
			\caption{Table 1}
		\end{table}
	
	\textit{Perform the following tasks:}
	
	Run a complete analysis on the data from Table 1. Assume the two-way factorial design. For full credit:
	
	\begin{enumerate}
		
		\item identify the research question
		\item State and estimate the statistical model
		\item Include exploratory plots to investigate the research question and interaction
		\item Include exploratory plots to assess the model assumptions.
		\item Test the interaction and main effects.
		\item Study the interaction in more detail using exploratory plots and by testing the relevant simple effects.
	\end{enumerate}
	
	\textbf{Answer}
	
	
	
	
	\clearpage
	
	\section{Question 2}
	
	{\color{deepred}\fontspec{Georgia}  ANCOVA with Interaction}
	
	Suppose that $(Y)$ is some continuous biological measurement and we assume the non-additive ANCOVA model:
	
	\begin{equation}
		Y_i = \beta_0 + \beta_1 x_{i1} + \beta_2 x_{i2} + \beta_3 x_{i1} x_{i,2} + \epsilon_i, \quad i = 1, \ldots, n, \quad \epsilon_i \overset{iid}{\sim} N(0, \sigma^2),
	\end{equation}
	
	where
	
	\[
	x_{i1} = \begin{cases} 1 & \text{if drug group} \\ 0 & \text{if control group,} \end{cases} \qquad x_{i2} = \text{continuous covariate.}
	\]
	
	\textit{Set-up}
	
	\begin{itemize}
		\item Further, suppose that the first $n_c$ cases are in the control group ($x_{i1} = 0$) and the following $n_D$ cases are in the drug group ($x_{1i} = 1$).
		
		\item Let $\bar{y}_C$, $\bar{x}_C$, $\bar{y}_D$ and $\bar{x}_D$ denote the respective group means of the response $Y$ and continuous covariate $x_2$ for the respective control and drug groups.
		
		\item Also consider two SLR models corresponding to the control group and drug group, each regressing $Y$ on continuous covariate $x_2$:
		
		\[
		Y_i = \beta_{0C} + \beta_C x_{i2} + \epsilon_i, \quad i = 1, \ldots, n_C, \quad \epsilon_i \overset{iid}{\sim} N(0, \sigma^2)
		\]
		
		\[
		Y_i = \beta_{0D} + \beta_D x_{i2} + \epsilon_i, \quad i = n_C + 1, \ldots, n, \quad \epsilon_i \overset{iid}{\sim} N(0, \sigma^2)
		\]
		
		Denote the estimated slopes of the above SLR models by $\hat{\beta}_C$ and $\hat{\beta}_D$.
		
		\item This exercise is similar to the \textit{``Linear Regression Approach to Single Factor ANCOVA''} Section from the lecture notes. Collectively, you will solve this problem in terms of the symbols $\bar{y}_C$, $\bar{x}_C$, $\bar{y}_D$, $\bar{x}_D$, $\hat{\beta}_C$ and $\hat{\beta}_D$.
	\end{itemize}

	\textit{Perform the following tasks:}

\begin{enumerate}
	\item Write an expression for all of the estimated coefficients of the non-additive linear regression model (1) in terms of the symbols $\bar{y}_C$, $\bar{x}_C$, $\bar{y}_D$, $\bar{x}_D$, $\hat{\beta}_C$ and $\hat{\beta}_D$. More specifically, identify the least squares estimators $\hat{\beta}_0$, $\hat{\beta}_1$, $\hat{\beta}_2$ and $\hat{\beta}_3$ of the estimated model
	\[
	Y_i = \hat{\beta}_0 + \hat{\beta}_1 x_{i1} + \hat{\beta}_2 x_{i2} + \hat{\beta}_3 x_{i1} x_{i,2}.
	\]
	In principle, you can solve a $(4 \times 4)$ inverse but that is not recommended. You should be able to solve this problem in a more intuitive fashion. The solution should be on the next two pages.
	
	\item Now consider centering the continuous covariate $x_2$, i.e., define the model
	\begin{equation}
		Y_i = \beta_0 + \beta_1 x_{i1} + \beta_2 u_{i2} + \beta_3 x_{i1} u_{x,i2} + \epsilon_i, \quad i = 1, \ldots, n, \quad \epsilon_i \overset{iid}{\sim} N(0, \sigma^2),
	\end{equation}
	where
	\[
	x_{i1} = \begin{cases} 1 & \text{if drug group} \\ 0 & \text{if control group,} \end{cases} \qquad u_{i2} = x_{i2} - \bar{x},
	\]
	and $\bar{x}$ is the sample average of continuous covariate $x_2$. Compute the least squares estimated coefficient $\hat{\beta}_1$ based on the above model (2).
	
	\item Interpret the estimated coefficient $\hat{\beta}_1$ based on both models (1) and (2). The expressions derived in problem parts 2.i and 2.ii should help with your interpretations.
\end{enumerate}
	
	\textbf{Answer}
	
	\clearpage
	
	\section{Question 3}
	
	{\color{deepred}\fontspec{Georgia}  Two-Way Model OLS}
	
	Consider the following equivalent two-factor model, which gives the mean of each $(i, j)^{th}$ cell:
	
	\begin{equation}
		Y_{ijk} = \mu_{ij} + \epsilon_{ijk}, \quad \epsilon_{ijk} \overset{iid}{\sim} N(0, \sigma^2),
	\end{equation}
	
	\[
	j = 1, \ldots, J, \quad i = 1, \ldots, I, \quad k = 1, \ldots, K,
	\]
	
	Formulate the least squares objective and derive the least squares estimators of $\mu_{ij}$.
	
	\textbf{Answer}
	
	\clearpage
	
	\section{Question 4}
	
	{\color{deepred}\fontspec{Georgia} Expected mean squares for Random Effects Model }
	
	Assume the one-way random effects ANOVA model:
	
	\begin{equation}
		Y_{ij} = \mu_i + \epsilon_{ij}, \quad \mu_i \overset{iid}{\sim} N(\mu, \sigma^2_A), \quad \epsilon_{ij} \overset{iid}{\sim} N(0, \sigma^2_e), \quad i = 1, \ldots, I,
	\end{equation}
	
	where $\mu_i$ are independent of $\epsilon_{ij}$.
	
	Under the above model statement, prove the following identities:
	
	\[
	E[MSA] = E[SSA/(I-1)] = J\sigma^2_A + \sigma^2_e
	\]
	
	\[
	E[MSE] = E[SSE/(IJ-I)] = \sigma^2_e
	\]
	
	\textbf{Answer}
	
	\clearpage
	
	\section{Question 5}
	
	{\color{deepred}\fontspec{Georgia}  Two-way-Table and Logistic Regression}
	
	Consider a study investigating the association between marijuana usage in college students and parental usage of drugs and alcohol. The following two-way frequency table summarizes the results. Each of 445 college students was classified according to both frequency of marijuana use and parental use of alcohol and psychoactive drugs.
	
	\begin{table}[h]
		\centering
		\begin{tabular}{llcc}
			\hline
			& & \multicolumn{2}{c}{Marijuana Use} \\
			& & No & Yes \\
			\hline
			Parental Use of & No & 141 & 94 \\
			Alcohol and Drugs & Yes & 85 & 125 \\
			\hline
		\end{tabular}
	\end{table}
	
	Use R to test if the \textit{odds ratio} of a college student that uses marijuana who came from a household without parental drug \& alcohol use verses a household with parental drug \& alcohol use statistically differs from 1. I.e., at $5\%$ significance, test if the odds ratio $\Theta = e^{\beta_1}$ statistically differs from 1.

	
	\textbf{Answer}
	
	\clearpage
	
	\section{Question 6}
	
	{\color{deepred}\fontspec{Georgia} Logistic Regression Application}
	
	A local health clinic sent fliers to its clients to encourage everyone, but especially older persons at high risk of complications, to get a flu shot in time for protection against an expected flu epidemic. In a pilot follow up study, 159 clients were randomly selected and asked whether they actually received a flu shot. A client who received a flu shot was coded $Y = 1$ and a client who did not receive a flu shot was coded $Y = 0$. In addition, data were collected on their age ($X_1$) and their health awareness. The latter data were combined into a health awareness index ($X_2$), for which higher values indicate greater awareness. Also included in the data was client gender, where males were coded $X_3 = 1$ and females were coded $X_3 = 0$. The data set \texttt{FluShots.txt} is Posted on Canvas.
	
	\textit{Perform the following tasks:}
	
	\begin{enumerate}
		\item Find the maximum likelihood estimators of $\beta_0, \beta_1, \beta_2, \beta_3$. State the fitted response function.
		\item Obtain $\exp(\hat{\beta}_1)$, $\exp(\hat{\beta}_2)$, $\exp(\hat{\beta}_3)$. Interpret these numbers.
		\item What is the estimated probability that male clients aged 55 with a health awareness index of 60 will receive a flu shot? Estimate this probability with $95\%$ confidence.
		\item Construct a ROC plot for the above model. Interpret the ROC plot.
	\end{enumerate}

	\textbf{Answer}
	
	\clearpage
	
	\section{Question 7}
	
	{\color{deepred}\fontspec{Georgia} Iteratively Re-weighted Least Squares}
	
	Consider the simple logistic regression model:
	
	\[
	Y_1, \ldots, Y_n \overset{ind}{\sim} Bernoulli(p_i),
	\]
	
	where
	
	\[
	p_i = \frac{e^{(\beta_0 + \beta_1 x_i)}}{1 + e^{(\beta_0 + \beta_1 x_i)}}, \quad i = 1, 2, \ldots, n.
	\]
	
	Further denote $\mathbf{y}$ as the response vector
	
	\[
	\mathbf{y} = [y_1 \ y_2 \ \cdots \ y_n]^T,
	\]
	
	$\mathbf{p}$ as the probability vector
	
	\[
	\mathbf{p} = [p_1 \ p_2 \ \cdots \ p_n]^T,
	\]
	
	$\mathbf{W}$ as the diagonal matrix of variances
	
	\[
	\mathbf{W} = \text{diag}\{p_1(1-p_1), p_2(1-p_2), \ldots, p_n(1-p_n)\},
	\]
	
	and $\mathbf{X}$ as the design matrix:
	
	\[
	\mathbf{X} = \begin{pmatrix} 1 & x_1 \\ 1 & x_2 \\ \vdots & \vdots \\ 1 & x_n \end{pmatrix}.
	\]
	
	\textit{Perform the following tasks:}
	
	\begin{enumerate}
		\item Set up the log-likelihood for the simple logistic regression model, i.e., set up
		\[
		\ell(\beta_0, \beta_1) = \ell(\beta_0, \beta_1; y_1, \ldots, y_n) = \log\{\mathcal{L}(\beta_0, \beta_1; y_1, \ldots, y_n)\}
		\]
		
		\item Show that the gradient of the log-likelihood can be written as:
		\[
		\nabla[\ell(\beta_0, \beta_1)] = \mathbf{X}^T(\mathbf{y} - \mathbf{p})
		\]
		
		\item Show that the Hessian of the log-likelihood can be written as:
		\[
		H_\ell(\beta_0, \beta_1) = -\mathbf{X}^T\mathbf{W}\mathbf{X}
		\]
		
		\item The update function of Newton's method is:
		\[
		\beta_t = \beta_{(t-1)} - [H_\ell(\beta)]^{-1}\nabla[\ell(\beta)]
		\]
		Newton's method searches for the extremum (min or max). Thus the update function of Newton's method for optimizing the logistic regression log-likelihood (or negative log-likelihood) is:
		\[
		\beta_t = \beta_{(t-1)} + [\mathbf{X}^T\mathbf{W}_{(t-1)}\mathbf{X}]^{-1}\mathbf{X}^T(\mathbf{y} - \mathbf{p}_{(t-1)})
		\]
		Manually code up Newton's method and test your algorithm on the same dataset from Problem 5. Check that your manual Newton's method algorithm produces the same result as the built-in function from R or \texttt{Python}. Run your Newton's method algorithm for enough iterations that the estimates stabilize.
	\end{enumerate}
	
	\textbf{Answer}
	
	\clearpage
	
	\section{Question 8}
	{\color{deepred}\fontspec{Georgia}  Mixed Model Estimation and the REML Objective}
		
		A clinical researcher wishes to evaluate the effectiveness of three drug treatments for reducing blood pressure. The treatments under consideration are \textit{Control}, \textit{Low Dose}, and \textit{High Dose}. To ensure that the results are generalizable, the study is conducted across several hospitals.
		
		Four hospitals are randomly selected from a district of 100 hospitals. Within each hospital, five patients are randomly assigned to each of the three treatments. At the end of the study period, the reduction in blood pressure is recorded for each patient.
		
		Consider the balanced \textit{two-factor mixed effects ANOVA model}
		
		\[
		Y_{ijk} = \mu + \alpha_i + \beta_j + \epsilon_{ijk}, \quad i = 1, 2, 3, \quad j = 1, 2, 3, 4, \quad k = 1, \ldots, 5,
		\]
		
		where
		
		\[
		\sum_{i=1}^{3} \alpha_i = 0, \quad \beta_j \overset{iid}{\sim} N(0, \sigma^2_B), \quad \epsilon_{ijk} \overset{iid}{\sim} N(0, \sigma^2),
		\]
		
		and the random effects $\beta_j$ are independent of the errors $\epsilon_{ijk}$.
		
		Formulate this same model using the {General Linear Mixed Model} framework:
		
		\[
		\mathbf{Y} = \mathbf{X}\boldsymbol{\beta} + \mathbf{Z}\mathbf{u} + \boldsymbol{\epsilon},
		\]
		
		where
		
		\[
		\mathbf{u} \sim N(\mathbf{0}, \mathbf{G}), \quad \boldsymbol{\epsilon} \sim N(\mathbf{0}, \boldsymbol{\Sigma}).
		\]
		
		In this problem, note that
		
		\[
		\mathbf{G} = \sigma^2_B \mathbf{I}_4, \quad \boldsymbol{\Sigma} = \sigma^2 \mathbf{I}_n,
		\]
		
		where $n = 3 \cdot 4 \cdot 5 = 60$. To align with the R REML function \texttt{lme()}, we use the \textit{linear regression parameterization} for the fixed effects. That is,
		
		\[
		\mathbf{X}\boldsymbol{\beta} \longleftrightarrow \beta_0 + \beta_1 x_{i1} + \beta_2 x_{i2}, \quad i = 1, 2, \ldots, 60,
		\]
		
		where the indicator variables are defined as
		
		\[
		x_{i1} = \begin{cases} 1, & \text{if observation } i \text{ receives the Low Dose,} \\ 0, & \text{otherwise,} \end{cases} \qquad x_{i2} = \begin{cases} 1, & \text{if observation } i \text{ receives the High Dose,} \\ 0, & \text{otherwise.} \end{cases}
		\]
		
		Also recall that the REML log-likelihood is given by
		
		\[
		\ell_R(\boldsymbol{\theta}) = -\frac{1}{2}\left[\log|\mathbf{V}| + \log|\mathbf{X}^\top \mathbf{V}^{-1}\mathbf{X}| + (\mathbf{y} - \mathbf{X}\hat{\boldsymbol{\beta}})^\top \mathbf{V}^{-1}(\mathbf{y} - \mathbf{X}\hat{\boldsymbol{\beta}})\right],
		\]
		
		where
		
		\[
		\mathbf{V} = \mathbf{Z}\mathbf{G}\mathbf{Z}^\top + \boldsymbol{\Sigma},
		\]
		
		and
		
		\[
		\hat{\boldsymbol{\beta}} = (\mathbf{X}^\top \mathbf{V}^{-1} \mathbf{X})^{-1} \mathbf{X}^\top \mathbf{V}^{-1} \mathbf{y}.
		\]
		
		\textit{Perform the following tasks:}
		
		\begin{enumerate}
			\item Specify the design matrices $\mathbf{X}$ and $\mathbf{Z}$ by writing down their skeleton forms, and state the dimensions of each matrix.
			
			\item Define the REML objective function
			\[
			Q(\boldsymbol{\eta}) = -\ell_R(\boldsymbol{\theta}),
			\]
			where the variance parameters are written on the log scale as
			\[
			\boldsymbol{\eta}^\top = \begin{pmatrix} \eta_1 & \eta_2 \end{pmatrix}, \quad \boldsymbol{\theta}^\top = \begin{pmatrix} e^{\eta_1} & e^{\eta_2} \end{pmatrix}.
			\]
			Interpret $e^{\eta_1}$ and $e^{\eta_2}$, and evaluate $Q(\boldsymbol{\eta})$ at
			\[
			\boldsymbol{\eta}^\top = \begin{pmatrix} 1 & 1 \end{pmatrix}.
			\]
			
			\item Use a Newton-based or gradient-based optimization routine in R or \texttt{Python} to minimize the objective function $Q(\boldsymbol{\eta})$. That is, solve for
			\[
			\hat{\boldsymbol{\eta}} = \arg\min_{\boldsymbol{\eta}} Q(\boldsymbol{\eta}).
			\]
			Then compute the corresponding variance component estimates
			\[
			\hat{\boldsymbol{\theta}}^\top = \begin{pmatrix} e^{\hat{\eta}_1} & e^{\hat{\eta}_2} \end{pmatrix}.
			\]
			
			\item Compute the estimated fixed-effects coefficients using the variance component estimates obtained in Problem 8.iii. That is, compute
			\[
			\hat{\boldsymbol{\beta}} = (\mathbf{X}^\top \hat{\mathbf{V}}^{-1} \mathbf{X})^{-1} \mathbf{X}^\top \hat{\mathbf{V}}^{-1} \mathbf{y}.
			\]
			
			\item Construct and carry out contrasts to test the following hypotheses:
			\begin{enumerate}
				\item Does the respondents' blood pressure reduction differ statistically between the High Dose and the control group?
				\item Does the respondents' blood pressure reduction differ statistically between the High Dose and the Low Dose group?
				\item Does respondents' blood pressure reduction differ statistically between the control group and both the High Dose and Low Dose groups?
			\end{enumerate}
		\end{enumerate}
		\textbf{Answer}
	
	\clearpage
	
	\section{Question 9}
	{\color{deepred}\fontspec{Georgia}   Mixed Model Estimation Continued (Practice)}
		
		Consider the \textit{two-way mixed effects ANOVA model} given by
		
		\begin{equation}
			Y_{ijk} = \mu + \alpha_i + \beta_j + \gamma_{ij} + \epsilon_{ijk}, \quad i = 1, \ldots, I, \ j = 1, \ldots, J, \ k = 1, \ldots, K.
		\end{equation}
		
		\textit{Model Assumptions}
		
		The parameter $\mu$ represents an overall mean common to all observations. The effects $\alpha_i$ are fixed effects associated with factor $A$ and satisfy the identifiability constraint
		
		\[
		\sum_{i=1}^{I} \alpha_i = 0.
		\]
		
		The effects $\beta_j$ are random effects associated with factor $B$, assumed to be independently and identically distributed as
		
		\[
		\beta_j \overset{iid}{\sim} N(0, \sigma^2_B).
		\]
		
		The interaction terms $\gamma_{ij}$ represent random effects between factors $A$ and $B$, with
		
		\[
		\gamma_{ij} \overset{iid}{\sim} N\left(0, \frac{I-1}{I}\sigma^2_{AB}\right), \quad \sum_{i=1}^{I} \gamma_{ij} = 0 \text{ for all } j.
		\]
		
		The error terms $\epsilon_{ijk}$ are assumed to be independently and identically distributed as
		
		\[
		\epsilon_{ijk} \overset{iid}{\sim} N(0, \sigma^2).
		\]
		
		Finally, all random components $\{\beta_j\}$, $\{\gamma_{ij}\}$, and $\{\epsilon_{ijk}\}$ are assumed to be mutually independent.
		
		\textit{Replicate Problem 8 using the full two-way mixed effects model. This is a practice problem, highly recommended to support your learning, but it will not be graded.}
		
		
	\textbf{Answer}
	
	\clearpage
	
		
		
		
	\end{document}
